% Generated draft: requires mathematical and source review.
\documentclass[11pt]{article}
\usepackage[margin=25mm]{geometry}
\usepackage{fontspec}
\setmainfont{Latin Modern Roman}
\usepackage{amsmath,amssymb,mathtools,unicode-math}
\setmathfont{Latin Modern Math}
\usepackage{enumitem,needspace,xurl,hyperref}
\setlength{\parindent}{0pt}
\setlength{\parskip}{5pt}
\setlength{\emergencystretch}{3em}
\newcommand{\N}{\mathbb{N}}
\newcommand{\Z}{\mathbb{Z}}
\newcommand{\Q}{\mathbb{Q}}
\newcommand{\R}{\mathbb{R}}
\newcommand{\C}{\mathbb{C}}
\newcommand{\F}{\mathbb{F}}
\newcommand{\A}{\mathbb{A}}
\newcommand{\PP}{\mathbb{P}}
\newcommand{\E}{\mathbb{E}}
\newcommand{\eps}{\varepsilon}
\newcommand{\dd}{\,\mathrm{d}}
\begin{document}
\section*{UnsolvedMath ID: 11100008}
\textbf{Source problem code:} catalogue entry
\subsection{Definitions and mathematical statement}
For an integer stem \(j\), the stable homotopy group of spheres is
\[
 \pi_j^S=\varinjlim_r\pi_{j+r}(S^r),
\]
with suspension transition maps. Unstable groups retain the actual sphere dimension \(k\), as \(\pi_n(S^k)\). Fixing \(k\) and letting \(n\to\infty\) is not the stable colimit.

\emph{Major problem 8 (Mark Hovey's list).} Say something general about stable or unstable homotopy groups of spheres. As a suggested growth question, the source attributes to Ravenel the possibility that the ``size'' of the \(n\)-th homotopy group of \(S^k\) grows polynomially in \(n\), perhaps even cubically. Hovey explicitly says he presumes the suggestion meant stable homotopy, while also proposing an unstable question. The ambiguity between these regimes must remain explicit.

Once a numerical size function \(s\) and indexing convention are chosen, an eventual polynomial upper bound would mean
\[
 \exists C>0\ \exists d\geq0\ \exists n_0\ \forall n\geq n_0,
 \qquad s(n)\leq Cn^d,
\]
with the suggested cubic version using \(d=3\). This only explains a possible upper-bound meaning of the source's growth wording, not a supplied exact conjecture. Cardinality, logarithm of cardinality, minimal generator count, and prime-primary rank are different functions. Some unstable groups have infinite cardinality, and \(\pi_0^S\cong\Z\) is also infinite; a cardinality formulation must state its degree range. Positive stable stems are finite, so their ordinary integral rank is zero: the summary's repeated word ``rank'' cannot by itself define the intended nontrivial size.

For a fixed-sphere unstable version the bound's constants may depend on \(k\); the excerpt does not state otherwise or fix prime localization. No particular size convention, uniformity in \(k\), or asymptotic equivalent is silently chosen. The broad request for general information is retained rather than equated solely with the illustrative polynomial suggestion.
\subsection{Short English statement}
Seek general information on sphere homotopy groups, including Ravenel's suggested polynomial, possibly cubic, growth question. Stable versus fixed-sphere unstable indexing and the meaning of size are unresolved; ordinary rank, torsion exponent, and group order must not be conflated.
\subsection{Sources}
\begin{itemize}
\item Local UnsolvedMath registry, permanent ID 11100008, source code catalogue entry; proposer Mark Hovey, ``Major problems,'' item 8. Ravenel is credited with the growth suggestion, and Hovey expressly qualifies its intended stable interpretation.
\item \emph{Hovey's Algebraic Topology Problem List}, surviving mirror: \url{https://www-users.cse.umn.edu/~tlawson/hovey/}. Recorded access 2026-07-29; extraction html-ordered-list. The mirror has 65 top-level items in nine category pages.
\item The summary mentions Ravenel, chromatic-height and spectral-sequence estimates, and Cohen--Moore--Neisendorfer exponent bounds, cross-referring to AMR-110-0061. No full bibliography or precise growth theorem is supplied; its final recent-computation discussion is truncated.
\item Supplied status partially\_solved and PARTIAL-PROGRESS. The summary describes the polynomial/cubic question as unresolved but uses an undefined rank notion and supplies no independent audit. Current primary-literature status and redistribution terms remain NEEDS\_REVIEW. The registry assigns L4 whereas background records default L3.
\end{itemize}
\end{document}
